\chapter{Представительный запрет семейного
\texorpdfstring{$SU(2)_F$}{SU(2)F}-подъёма на
\texorpdfstring{$H_{15}/M_{35}$}{H15/M35}}

Предыдущая глава выбрала наиболее экономный кандидат продолжения:
попытаться прочитать уже построенный семейный $SO(3)$-перенос как
присоединённую тень одного локального $SU(2)_F$-переноса. Настоящий гейт
выполняет объявленный ранний стоп-тест. Монопольный профиль, потенциал и
новая константа связи не вводятся.

\section{Антикруговая сверка и точный носитель}

Тест опирается только на уже замороженные данные проекта. На физической
частичной половине наблюдаемый пакет имеет вид
\begin{equation}
 \mathcal H_{\rm light}=\mathbb C^3_{\rm fam}\otimes H_{15},
 \qquad \dim\mathcal H_{\rm light}=45,
\end{equation}
где
\begin{equation}
 H_{15}=Q_L^{(6)}\oplus L_L^{(2)}\oplus
 u_R^{(3)}\oplus d_R^{(3)}\oplus e_R^{(1)}.
 \label{eq:v5-su2-h15}
\end{equation}
В родительской семейной цепи частичная половина имеет размерность
$4+3+3=10$ и под непрерывными семейными вращениями раскладывается как
\begin{equation}
 \mathbb C^{10}_{F}\simeq\mathbf 1\oplus\mathbf 3\oplus
 \mathbf 3\oplus\mathbf 3.
 \label{eq:v5-su2-family-decomposition}
\end{equation}
KO6 добавляет сопряжённую половину и даёт $20$, а наблюдаемый множитель
$H_{15}$ приводит к $20\cdot15=300$. Связывающая алгебра остаётся
$M_{35}$ с углами $M_{20}$ и $M_{15}$.

Следовательно, задача состоит не в поиске произвольного представления
$SU(2)$ той же размерности, а в проверке, содержит ли именно этот модуль
нетривиальное полуцелое представление, не меняющее его физический смысл.

\section{Канонический подъём триплета не хранит центральный знак}

Существующий семейный триплет является векторным представлением
$SO(3)$. Его канонический подъём вдоль накрытия
\begin{equation}
 SU(2)\longrightarrow SO(3)
\end{equation}
есть представление спина $j=1$. Центральный элемент $-1\in SU(2)$
соответствует повороту на $2\pi$ и действует как
\begin{equation}
 \rho_{j=1}(-1)=+I_3.
 \label{eq:v5-su2-spin-one-center}
\end{equation}
Для фундаментального дублета спина $j=1/2$, напротив,
\begin{equation}
 \rho_{j=1/2}(-1)=-I_2.
 \label{eq:v5-su2-spin-half-center}
\end{equation}

Все слагаемые в \eqref{eq:v5-su2-family-decomposition} имеют целый спин.
Поэтому центр действует тривиально на семейных модулях размерностей $10$,
$20$, $45$ и $300$. Присоединённое $SO(3)$-чтение воспроизводится, но
потерянный проектором знак не восстанавливается.

\begin{theorem}[Центральное препятствие текущего семейного модуля]
Канонический $SU(2)$-подъём существующего семейного
$\mathbf 1\oplus3\mathbf 3$ содержит только целоспиновые представления.
Он факторизуется через $SO(3)$, и центральный элемент $-1$ действует
тождественно. Поэтому один текущий семейный модуль не может одновременно
служить векторным носителем проектора $P=nn^T$ и спинорным носителем
массы $n\cdot\boldsymbol\sigma$.
\end{theorem}

\section{Почему KO6-удвоение не является скрытым дублетом}

Естественно попытаться использовать пару ``частица--сопряжённая частица''
как фундаментальную $\mathbb C^2$. На этом множителе градуировка имеет вид
\begin{equation}
 \gamma_{\rm KO6}=\sigma_z\otimes I_{10}.
\end{equation}
Полное действие $SU(2)$ потребовало бы генераторов
$\sigma_x,\sigma_y,\sigma_z$. Но
\begin{equation}
 [\sigma_x\otimes I_{10},\gamma_{\rm KO6}]\ne0,
 \qquad
 [\sigma_y\otimes I_{10},\gamma_{\rm KO6}]\ne0,
\end{equation}
тогда как с градуировкой коммутирует только диагональное направление
$\sigma_z$. Таким образом, полный $SU(2)$ смешивает половины с
противоположной градуировкой, а также частичные и сопряжённые
стандарт-модельные представления. Сохранить без изменения можно лишь
диагональную $U(1)$, не фундаментальный семейный дублет.

\section{Почему слабый дублет не решает семейную задачу}

В $H_{15}$ уже существует фундаментальное действие слабой группы
$SU(2)_L$. Однако её центр действует знаком минус только на восьми левых
состояниях: трёх цветных дублетах $Q_L$ и одном дублете $L_L$. На семи
правых синглетах он действует тривиально. Поэтому центральное действие на
$H_{15}$ имеет спектр
\begin{equation}
 \operatorname{Spec}\rho_L(-1)=
 \{(-1)^{\times8},(+1)^{\times7}\}.
\end{equation}
Это не единый семейный знак и не действие на трёх поколениях. Отождествить
$SU(2)_F$ с $SU(2)_L$ означало бы изменить стандарт-модельные квантовые
числа и смешать семейную ось со слабым изоспином.

\section{Разложение
\texorpdfstring{$\mathbf3\to\mathbf2\oplus\mathbf1$}{3 в 2 плюс 1}}

Можно сохранить размерность семейного пространства, заменив триплет на
\begin{equation}
 \mathbb C^3_{\rm fam}\longmapsto\mathbf2\oplus\mathbf1.
 \label{eq:v5-su2-two-plus-one}
\end{equation}
Тогда центр имеет собственные значения $(-1,-1,+1)$ и действительно
появляется полуцелый сектор. Но это уже другое семейное представление:
исходный неприводимый $SO(3)$-триплет, тетраэдрическая геометрия и
канонический проектор одной инвариантной линии больше не сохраняются.

Если дополнительно считать этот $SU(2)_F$ динамической хиральной
калибровочной группой, на физической половине возникает по одному
семейному дублету для каждого из пятнадцати Weyl-каналов $H_{15}$. Число
дублетов равно $15$ и нечётно. Это попадает под глобальную аномалию
$SU(2)$ Виттена. KO6-сопряжённая половина не устраняет проблему
физической хиральной теории автоматически: для такого спасения пришлось
бы заново объявить вектороподобный спектр и повторить физическое чтение.

\section{Ориентационный дублет локального переноса}

В локальном квантовом блуждании уже появлялся множитель $\mathbb C^2$ для
двух встречных направлений распространения. Он действительно может нести
матрицы Паули и кинематически выглядит как нужный полуцелый модуль.
Однако он был добавлен после
\begin{equation}
 E=M_{20\times15}(\mathbb C)
\end{equation}
и не является частью $H_{15}$, $M_{35}$ или их замороженного следа.

Отождествление этого дублета с семейным спинором имеет только два чтения:
\begin{enumerate}
 \item добавить внешнее ассоциированное расслоение, то есть нарушить
 текущий бюджет модулей;
 \item связать направление распространения и семейный $SU(2)$ общей
 структурой типа $\operatorname{Spin}^h$.
\end{enumerate}
Второй вариант концептуально интересен, но является новой архитектурой:
он меняет правило симметрии, требует повторного аудита KO6, аномалий,
следа и феноменологии. Он не является прохождением текущего гейта.

\section{След и бюджет состояний}

Нормированный след $M_{35}$ фиксирует относительные веса существующих
углов $M_{20}$ и $M_{15}$. Он не нормирует внешний фундаментальный
дублет. Добавление полуцелого ассоциированного модуля требует либо нового
матричного угла, либо дополнительного тензорного множителя и тем самым
повторного вывода общей меры. Это именно тот скрытый секторный вес, который
стоп-критерий запрещал вводить молча.

\begin{nogocriterion}[Представительный запрет $H_{15}/M_{35}$]
На текущем носителе нельзя одновременно получить присоединённый семейный
$SO(3)$-триплет и нетривиальное центральное спинорное действие $-1$.
KO6-удвоение нарушает градуировку при полном $SU(2)$-действии; слабый
дублет имеет иные квантовые числа; замена $3\to2\oplus1$ разрушает
семейную геометрию и в хиральном калибровочном чтении даёт нечётное число
дублетов; ориентационный дублет лежит вне $H_{15}/M_{35}$.
\end{nogocriterion}

\section{Вердикт и условие переоткрытия}

Первый обязательный гейт основного маршрута не пройден. Это ранний и
полезный отрицательный результат: монопольную динамику, потенциал и
солитонный профиль строить не следует. Существующий $SO(3)$-перенос
сохраняется как векторная геометрия, но не является тенью уже имеющегося
фундаментального семейного спинора.

Переоткрытие допустимо только как отдельно объявленная новая версия с
одним из двух явных изменений:
\begin{enumerate}
 \item новый полуцелый ассоциированный модуль с полным аудитом состояний,
 аномалий и следа;
 \item структура типа $\operatorname{Spin}^h$, связывающая существующий
 ориентационный дублет распространения с семейным переносом, с повторным
 выводом KO6 и физического чтения.
\end{enumerate}
Ни один из этих вариантов не принимается автоматически.

\begin{protocol}
\begin{enumerate}
 \item семейное присоединённое $SO(3)$-чтение: \textbf{сохраняется};
 \item центральный знак на текущем семейном модуле: \textbf{отсутствует};
 \item KO6 как скрытый фундаментальный дублет: \textbf{запрещён};
 \item слабый $SU(2)_L$ как семейный подъём: \textbf{несовместим};
 \item замена $3\to2\oplus1$: \textbf{меняет модель и аномальна в
 минимальном хиральном чтении};
 \item ориентационный дублет: \textbf{условный ресурс новой версии};
 \item нормировка внешнего дублета следом $M_{35}$: \textbf{не выведена};
 \item первый стоп-тест: \textbf{не пройден};
 \item монопольная динамика в версии V: \textbf{не открывается};
 \item физическое замыкание: \textbf{не достигнуто}.
\end{enumerate}
\end{protocol}

Машинный протокол сохранён в
\path{s2t_v5_su2_family_lift_h15_representation_gate_results.json}.